\LessonHeader{35}{Why Teenage Boys Stop Reading}{十几岁的男孩为什么不再读书}{教育新闻 · Education News}{话题\ 青少年阅读·调查报告\enspace|\enspace 新词\ \textbf{3}\enspace|\enspace 新义\ \textbf{18}\enspace|\enspace 短语\ \textbf{12}\enspace|\enspace 全文\ \textbf{822}\ 词}{Lesson 35\enspace·\enspace Why Teenage Boys Stop Reading}
\begin{lessonbody}
\PARA{1}How \RV{}{often}{} do you \RV{}{pick}{} up a book just \PH{for\ \RV{}{fun}{?}}{for fun}{为了消遣} If you are a boy in the \PN{}{UK}{}{/ˌjuː ˈkeɪ/}{n.}{英国} \DV{}{aged}{}{age + -ed → 形容词}{adj.}{……岁的} 14 to 16, the \RV{}{answer}{} is \RV{}{probably}{} “not very \AG{}{often}{”.} New \RV{}{research}{} has \RVB{}{found}{}{find} that \RVB{}{fewer}{}{few} \RV{}{than}{} one \RV{}{in}{} ten boys of that \RV{}{age}{} read \RV{}{every}{} day. \DV{}{Schoolwork}{,}{school + work → 名词}{n.[U]}{学校功课} \RVB{}{screens}{}{screen} and \RV{}{sport}{,} it says, are \NS{}{pushing}{}{新义}{push}{/pʊʃ/}{vt.}{挤（出）；排挤} reading \PH{for\ \RV{}{pleasure}{}}{for pleasure}{为了乐趣} out of \RV{}{teenage}{} \RV{}{life}{.}

\PARA{2}The \PN{}{National Literacy Trust}{,}{/ˌnæʃnəl ˈlɪtərəsi trʌst/}{n.}{英国国家读写素养基金会} a \NS{}{group}{}{新义}{}{/ɡruːp/}{n.[C]}{团体；组织} that helps children read and write, asked \RV{}{nearly}{} 115,000 children and \RV{}{young}{} people \AG{}{aged}{} five to 18 about their \NW{}{reading}{.}{/ˈriːdɪŋ/}{n.[U]}{阅读} \RV{}{Almost}{} \RV{}{half}{} (46.9\%) of children \AG{}{aged}{} eight to 11 said they \RVB{}{enjoyed}{}{enjoy} reading. \RV{}{Among}{} 11- to 14-year-\RVB{}{olds}{,}{old} that \RV{}{number}{} \RVB{}{fell}{}{fall} to \RV{}{less}{} \AG{}{than}{} a \RV{}{third}{} (29.5\%), and for 14- to 16-year-\AG{}{olds}{} it \AG{}{fell}{} \RV{}{even}{} \RV{}{further}{,} to 28.6\%.

\PARA{3}\RV{}{Daily}{} \AG{}{reading}{} \RVB{}{drops}{}{drop} \RV{}{in}{} the \RV{}{same}{} way. Of children \AG{}{aged}{} eight to 11, 31.1\% read \AG{}{every}{} day. Between 11 and 14, it is 17.1\%, and from 14 to 16, it is only 14\%. By the \RV{}{middle}{} of the \AG{}{teenage}{} years, the \NS{}{report}{}{新义}{}{/rɪˈpɔːt/}{n.[C]}{（调查）报告} says, reading \AG{}{every}{} day has \RV{}{become}{} \NW{}{unusual}{.}{/ʌnˈjuːʒuəl/}{adj.}{少见的}

\PARA{4}\AG{}{Reading}{} \AGB{}{falls}{}{fall} for \RV{}{both}{} boys and girls in the \RV{}{early}{} \AG{}{teenage}{} years. But in the \NS{}{later}{}{新词性}{}{/ˈleɪtə(r)/}{adj.}{后期的} \AG{}{teenage}{} years there are \RVB{“}{signs}{}{sign} of \DV{}{recovery}{”}{recover + -y → 名词}{n.[U]}{恢复} \AG{}{among}{} girls, \RV{}{while}{} the \AGB{}{numbers}{}{number} for boys \RV{}{stay}{} \RV{}{low}{.} Over a \AG{}{third}{} of girls (36\%) \AG{}{aged}{} eight to 11 \NS{}{reported}{}{新义}{report}{/rɪˈpɔːt/}{vt.}{报告说；称} reading \NS{}{daily}{,}{新词性}{}{/ˈdeɪli/}{adv.}{每天} \RVB{}{compared}{}{compare} with about a \RV{}{quarter}{} (26.3\%) of boys of the \AG{}{same}{} \AG{}{age}{.} By 14 to 16, this had \AGB{}{fallen}{}{fall} to 17.6\% of girls and just 9.8\% of boys.

\PARA{5}\SMK{y}\PN{}{Jonathan Douglas}{,}{/ˌdʒɒnəθən ˈdʌɡləs/}{n.}{乔纳森·道格拉斯} who \RVB{}{leads}{}{lead} the \NS{}{trust}{,}{新词性}{}{/trʌst/}{n.[C]}{基金会} said: “Over the \RV{}{past}{} 20 years, children and \AG{}{young}{} people’s \DV{}{enjoyment}{}{enjoy + -ment → 名词}{n.[U]}{喜爱} of \AG{}{reading}{} and their \AG{}{daily}{} \AG{}{reading}{} \RVB{}{habits}{}{habit} have \AGB{}{fallen}{}{fall} \GL{}{steeply}{}{/ˈstiːpli/}{adv.}{急剧地} to their \AGB{}{lowest}{}{low} \NS{}{levels}{}{新义}{level}{/ˈlevl/}{n.[C]}{水平} \PH{on\ \RV{}{record}{,}}{on record}{有记录以来} \RV{}{with}{} the \NS{}{drop}{}{新义}{}{/drɒp/}{n.[C]}{下降} \RV{}{especially}{} \GL{}{pronounced}{}{/prəˈnaʊnst/}{adj.}{明显的} for \RVB{}{teenagers}{.”}{teenager}\EMK{y}{1} 
\YI{1}{负责该基金会的乔纳森·道格拉斯说：“过去20年里，儿童和青少年对阅读的喜爱和每天阅读的习惯都急剧下降，降到了有记录以来的最低水平，青少年下降得尤其明显。”}
\RV{}{As}{} \AG{}{young}{} people \RV{}{grow}{} older, he \RVB{}{explained}{,}{explain} \AG{}{reading}{} has to \RV{}{fight}{} for their time \RV{}{against}{} many other things: more \RV{}{pressure}{,} \RVB{}{changing}{}{change} \RVB{}{routines}{,}{routine} new \NS{}{interests}{}{新义}{interest}{/ˈɪntrəst/}{n.[C]}{爱好} and \RVB{}{activities}{,}{activity} and more \DV{}{freedom}{.}{free + -dom → 名词}{n.[U]}{自由} \PH{\RV{“}{As}{}\ a\ \RV{}{result}{,}}{as a result}{因此} \AG{}{reading}{} can \AG{}{become}{} \RVB{}{easier}{}{easy} to \AG{}{push}{} out, \DV{}{particularly}{}{particular + -ly → 副词}{adv.}{尤其} when it is not \RV{}{part}{} of \AG{}{daily}{} \AG{}{life}{.”} \AG{}{Teenage}{} boys’ \AG{}{reading}{,} he \RVB{}{added}{,}{add} is \AG{}{especially}{} \AG{}{easy}{} to \RV{}{lose}{.}

\PARA{6}The \AG{}{report}{} looks \RV{}{closely}{} at \AGB{}{answers}{}{answer} from 80,000 \RVB{}{students}{}{student} \AG{}{aged}{} 11 to 16 who took \RV{}{part}{} in the \RV{}{study}{.} It comes out \RV{}{during}{} the \PN{}{National Year of Reading}{,}{/ˌnæʃnəl ˌjɪər əv ˈriːdɪŋ/}{n.}{全国阅读年} a year of work \RV{}{across}{} the UK to get more people, \AG{}{especially}{} children and \AG{}{young}{} people, reading for \AG{}{pleasure}{.}

\PARA{7}\PN{}{Phil Earle}{}{/ˌfɪl ˈɜːl/}{n.}{菲尔·厄尔} is a \RV{}{writer}{} \RV{}{whose}{} \RV{}{novel}{} \PN{}{When the Sky Falls}{}{/wen ðə ˈskaɪ fɔːlz/}{n.}{《天塌之时》} \RVB{}{won}{}{win} \PN{}{Children’s Book of the Year}{}{/ˌtʃɪldrənz ˌbʊk əv ðə ˈjɪə(r)/}{n.}{年度童书奖} at the \PN{}{British Book Awards}{}{/ˌbrɪtɪʃ ˈbʊk əˌwɔːdz/}{n.}{英国图书奖} in 2022. “If we are \NS{}{serious}{}{新义}{}{/ˈsɪəriəs/}{adj.}{严肃认真的} about this \RV{}{situation}{} then it won’t \RV{}{mean}{} a year of work, it \AGB{}{means}{}{mean} a \GL{}{generation}{}{/ˌdʒenəˈreɪʃn/}{n.[C]}{一代人} of it,” he said. \RV{}{Everyone}{,} from the \RV{}{government}{} and \RVB{}{teachers}{}{teacher} to \RV{}{TV}{,} \RV{}{radio}{,} \RV{}{social}{} \RV{}{media}{,} \DV{}{publishers}{}{publish + -er → 名词}{n.[C]}{出版商} and \AGB{}{writers}{,}{writer} will \RV{}{need}{} to help.

\PARA{8}It is time, he said, to \RV{}{end}{} \RVB{}{centuries}{}{century} of \GL{}{snobbery}{}{/ˈsnɒbəri/}{n.[U]}{势利；看不起} about \AG{}{reading}{.} \SMK{y}Children \RV{}{should}{} be our \NS{}{guides}{}{新义}{guide}{/ɡaɪd/}{n.[C]}{向导} \PH{when\ it\ comes\ to}{when it comes to}{说到} what they want to read, and children’s \AGB{}{writers}{}{writer} \AG{}{should}{} be seen more on \AG{}{TV}{} and \RV{}{online}{,} so that \RVB{}{parents}{,}{parent} and more \DV{}{importantly}{}{important + -ly → 副词}{adv.}{重要的是} \AG{}{young}{} \RVB{}{readers}{}{reader} \RV{}{themselves}{,} know who they are.\EMK{y}{2} 
\YI{2}{说到孩子想读什么，孩子应该给我们引路；儿童作家也应该更多地出现在电视和网上，让家长，更重要的是让年轻读者自己，知道他们是谁。}
“It won’t be \AG{}{easy}{,} and it won’t be \RV{}{quick}{,} but we can \RV{}{keep}{} our \AG{}{readers}{} by \RVB{}{reminding}{}{remind} them \NS{}{way}{}{新词性}{}{/weɪ/}{adv.}{远远（早于）} before they’re \AG{}{teenagers}{} that \AG{}{reading}{} isn’t just about \AGB{}{novels}{}{novel} too.”

\PARA{9}The \AG{}{study}{,} \RVB{}{shared}{}{share} first with the \PN{}{Guardian}{,}{/ˈɡɑːdiən/}{n.}{《卫报》} also \NS{}{includes}{}{新词性}{include}{/ɪnˈkluːd/}{vt.}{包括} \AG{}{almost}{} 50,000 \RVB{}{comments}{}{comment} from 11- to 16-year-\AG{}{olds}{} about how \AG{}{reading}{} \RVB{}{fits}{,}{fit} or does not \AG{}{fit}{,} into their \RV{}{everyday}{} \AGB{}{lives}{.}{life} Many \AG{}{teenage}{} boys had good things to say. They \RVB{}{found}{}{find} \AG{}{reading}{} \RV{}{calm}{} and \DV{}{peaceful}{,}{peace + -ful → 形容词}{adj.}{宁静的} and for many of them it is \RV{}{still}{} a \RV{}{real}{} \NS{}{pleasure}{.}{新义}{}{/ˈpleʒə(r)/}{n.[C]}{乐事} “If I’m \GL{}{upset}{,}{/ʌpˈset/}{adj.}{难过的} I can read a book to \PH{\NS{}{calm}{}{新词性}{}{/kɑːm/}{vi.}{平静下来}\ down,”}{calm down}{平静下来} said one boy. \AG{“}{Reading}{} is \RV{}{great}{} for when you have a bad day at school. You can just read a book to stop \PH{\RVB{}{worrying}{}{worry}\ about}{worry about}{为……担心} it,” said \RV{}{another}{.} Some also \AG{}{enjoyed}{} \RVB{}{learning}{}{learn} \RV{}{through}{} \AG{}{reading}{.} “I like reading because I \GL{}{gain}{}{/ɡeɪn/}{vt.}{获得} better \RV{}{knowledge}{,”} one said.

\PARA{10}\SMK{y}\AG{}{Often}{,} \RV{}{however}{,} reading for \AG{}{fun}{} \PH{\RVB{}{loses}{}{lose}\ out\ to}{lose out to}{输给} \AG{}{schoolwork}{.}\EMK{y}{3} 
\YI{3}{不过，为乐趣而读书常常要给学校功课让路。}
“I don’t have the time to \AG{}{enjoy}{} books when I am \RV{}{already}{} \GL{}{drowning}{}{/ˈdraʊnɪŋ/}{vi.}{淹没；被压垮} in \DV{}{textbooks}{,”}{text + book → 名词}{n.[C]}{课本} said one boy. Many \PH{\RV{}{would}{}\ \RV{}{rather}{}}{would rather}{宁愿} \PH{\RV{}{spend}{}\ their\ time\ on}{spend time on}{把时间花在} \AG{}{sport}{,} \RVB{}{games}{,}{game} and their \RV{}{family}{} and friends. Others liked reading on a \AG{}{screen}{.} “I don’t \AG{}{enjoy}{} reading a book, but I \AG{}{enjoy}{} reading things on my \RV{}{phone}{,} \RV{}{computer}{} and more,” one said.

\PARA{11}Some boys are \PH{put\ \RV{}{off}{}}{put off}{使失去兴趣} because \AG{}{reading}{} \RV{}{itself}{} is \RV{}{hard}{} for them. “I am \GL{}{dyslexic}{}{/dɪsˈleksɪk/}{adj.}{有阅读障碍的} and don’t \AG{}{enjoy}{} it \NS{}{as}{}{新义}{}{/æz/}{conj.}{因为} it’s too \AG{}{hard}{,} my \RV{}{attention}{} \GL{}{span}{}{/spæn/}{n.[C]}{持续时间} is too \RV{}{short}{} and I \AG{}{find}{} other things more \RV{}{fun}{,”} one boy \AG{}{explained}{.}

\PARA{12}So what helps? \AG{“}{Teenagers}{} \AG{}{themselves}{} told us they are more \RV{}{likely}{} to read when it \NS{}{connects}{}{新义}{connect}{/kəˈnekt/}{vi.}{与……有关联} to their \AG{}{interests}{,} \RV{}{such}{} \RV{}{as}{} \RV{}{football}{,} \RV{}{video}{} \AG{}{games}{} and \RV{}{friendship}{,”} Douglas said, and when they can \RV{}{choose}{} what and how they read, and when \AG{}{reading}{} \AG{}{fits}{} \RV{}{naturally}{} into their \AGB{}{lives}{.}{life} \AG{}{Young}{} people read \RV{}{text}{} \AG{}{in}{} many \RV{}{different}{} ways today, he \AG{}{added}{.} \SMK{y}Seeing how all of these can work \RV{}{together}{} with books will \PH{be\ \RV{}{key}{}\ to}{be key to}{是……的关键} helping more \AG{}{teenagers}{} \NS{}{develop}{}{新义}{}{/dɪˈveləp/}{vt.}{养成；形成} a \AG{}{reading}{} \AG{}{habit}{} that can help their \NW{}{learning}{,}{/ˈlɜːnɪŋ/}{n.[U]}{学习} their \DV{}{happiness}{}{happy + -ness → 名词}{n.[U]}{快乐} and their \RV{}{confidence}{} at an \RV{}{important}{} \RV{}{moment}{} in their \AGB{}{lives}{.}{life}\EMK{y}{4}
\YI{4}{看清这些方式怎样和书本一起发挥作用，将是帮助更多青少年养成阅读习惯的关键；在人生的重要时刻，这种习惯能帮助他们学习，让他们更快乐、更自信。}
\end{lessonbody}
\sourceline{Adapted from: “Only 10\% of boys aged 14-16 read daily for pleasure, National Literacy Trust finds”, The Guardian, 2026-02-22.}
\SelfTest{33}{\TW{reading}
\TW{unusual}
\TW{learning}
\TW{push\,\nstag{新义}}
\TW{group\,\nstag{新义}}
\TW{report\,\nstag{新义}}
\TW{later\,\nstag{新词性}}
\TW{daily\,\nstag{新词性}}
\TW{trust\,\nstag{新词性}}
\TW{level\,\nstag{新义}}
\TW{drop\,\nstag{新义}}
\TW{interest\,\nstag{新义}}
\TW{serious\,\nstag{新义}}
\TW{guide\,\nstag{新义}}
\TW{way\,\nstag{新词性}}
\TW{include\,\nstag{新词性}}
\TW{pleasure\,\nstag{新义}}
\TW{calm\,\nstag{新词性}}
\TW{as\,\nstag{新义}}
\TW{connect\,\nstag{新义}}
\TW{develop\,\nstag{新义}}
\TW{for fun}
\TW{for pleasure}
\TW{on record}
\TW{as a result}
\TW{when it comes to}
\TW{calm down}
\TW{worry about}
\TW{lose out to}
\TW{would rather}
\TW{spend time on}
\TW{put off}
\TW{be key to}}
\par\write-1{LESSONEND 35 \thepage}
